% GENERATED FILE. Edit canonical lessons, then run npm run generate:books.
% canonical-source-hash: fnv1a64:92376f391662f192
% canonical-lessons: GE-C413-schlafwagen, GE-C413-speisewagen, GE-C413-endstation, GE-C413-reservieren, GE-C413-weiterfahren, GE-R413-review-406-413

\chapter{Im Nachtzug}
\label{ch:ge-im-nachtzug}
\hlchaptermodality{413}

\begin{chapteropening}
\textbf{By the end of this chapter:} \emph{I can travel on a night train: the sleeping car, the dining car, the end of the line, booking and travelling on.}

Complete the last lesson of chapter 413: I can travel on a night train: the sleeping car, the dining car, the end of the line, booking and travelling on.
\end{chapteropening}

\section[der Schlafwagen]{der Schlafwagen — the sleeping car}

\label{lesson:GE-C413-schlafwagen}



\begin{quote}
Before the new one: say the German for the conductor, then the German for to check.
\end{quote}

\subsection*{You'll want to know: der Schlafwagen}
\textbf{der Schlafwagen} — \textquotedblleft{}the sleeping car\textquotedblright{}.

\begin{quote}
Im Schlafwagen schläft sie bis zum Morgen.
\end{quote}

\begin{quote}
\emph{In the sleeping car she sleeps until morning.}
\end{quote}

Same form in the plural. A cheaper car with couchettes is \emph{der Liegewagen}.

\begin{cousinweb}
\emph{Schlafwagen} is built from \emph{schlafen} 'to sleep' and \emph{Wagen} 'carriage', both inherited: \emph{Wagen} goes back to Proto-Germanic \emph{wagnaz} (EtymDB).
\end{cousinweb}

\begin{grammarlens}[title={What you've built}]
One more word for telling, explaining and describing at B1.
\end{grammarlens}

\subsection*{Your turn}
\begin{itemize}
\raggedright
  \item \emph{Say it:} \emph{der Schlafwagen}
  \item \emph{Say it:} \emph{der Schlafwagen}, once more
  \item \emph{Say it:} a sentence of your own with \emph{der Schlafwagen}
  \item \emph{Recall:} say the German for the conductor, then the German for to check, then say \emph{der Schlafwagen} again
\end{itemize}

\begin{tcolorbox}[breakable,colback=teal!4,colframe=teal!35!black,title={Before you move on}]
What does der Schlafwagen mean? (\textquotedblleft{}The sleeping car\textquotedblright{}.) Say it once more.
\end{tcolorbox}
\section[der Speisewagen]{der Speisewagen — the dining car, the restaurant car}

\label{lesson:GE-C413-speisewagen}



\begin{quote}
Before the new one: say the German for to check, then the German for the sleeping car.
\end{quote}

\subsection*{You'll want to know: der Speisewagen}
\textbf{der Speisewagen} — \textquotedblleft{}the dining car, the restaurant car\textquotedblright{}.

\begin{quote}
Im Speisewagen essen wir eine warme Suppe.
\end{quote}

\begin{quote}
\emph{In the dining car we eat a hot soup.}
\end{quote}

Same form in the plural. Also \emph{das Bordrestaurant}.

\begin{cousinweb}
\emph{Speisewagen} is built from \emph{Speise} 'food, dish', which goes back to Middle High German \emph{spīse} (EtymDB) and is usually traced to Latin \emph{expēnsa} 'outlay, expense' (Lewis \& Short), and \emph{Wagen} 'carriage'.
\end{cousinweb}

\begin{grammarlens}[title={What you've built}]
One more word for telling, explaining and describing at B1.
\end{grammarlens}

\subsection*{Your turn}
\begin{itemize}
\raggedright
  \item \emph{Say it:} \emph{der Speisewagen}
  \item \emph{Say it:} \emph{der Speisewagen}, once more
  \item \emph{Say it:} a sentence of your own with \emph{der Speisewagen}
  \item \emph{Recall:} say the German for to check, then the German for the sleeping car, then say \emph{der Speisewagen} again
\end{itemize}

\begin{tcolorbox}[breakable,colback=teal!4,colframe=teal!35!black,title={Before you move on}]
What does der Speisewagen mean? (\textquotedblleft{}The dining car, the restaurant car\textquotedblright{}.) Say it once more.
\end{tcolorbox}
\section[die Endstation]{die Endstation — the end of the line, the terminus}

\label{lesson:GE-C413-endstation}



\begin{quote}
Before the new one: say the German for the sleeping car, then the German for the dining car.
\end{quote}

\subsection*{You'll want to know: die Endstation}
\textbf{die Endstation} — \textquotedblleft{}the end of the line, the terminus\textquotedblright{}.

\begin{quote}
An der Endstation steigen alle aus.
\end{quote}

\begin{quote}
\emph{At the end of the line everyone gets off.}
\end{quote}

Plural: \emph{die Endstationen}. Also figurative: \emph{Endstation} — this is as far as it goes.

\begin{cousinweb}
\emph{Endstation} is built from \emph{End(e)} 'end', inherited from Proto-Germanic \emph{andiją}, and \emph{Station}, from Latin \emph{statiō} 'a standing, a post' (Lewis \& Short).
\end{cousinweb}

\begin{grammarlens}[title={What you've built}]
One more word for telling, explaining and describing at B1.
\end{grammarlens}

\subsection*{Your turn}
\begin{itemize}
\raggedright
  \item \emph{Say it:} \emph{die Endstation}
  \item \emph{Say it:} \emph{die Endstation}, once more
  \item \emph{Say it:} a sentence of your own with \emph{die Endstation}
  \item \emph{Recall:} say the German for the sleeping car, then the German for the dining car, then say \emph{die Endstation} again
\end{itemize}

\begin{tcolorbox}[breakable,colback=teal!4,colframe=teal!35!black,title={Before you move on}]
What does die Endstation mean? (\textquotedblleft{}The end of the line, the terminus\textquotedblright{}.) Say it once more.
\end{tcolorbox}
\section[reservieren]{reservieren — to book, to reserve}

\label{lesson:GE-C413-reservieren}



\begin{quote}
Before the new one: say the German for the dining car, then the German for the end of the line.
\end{quote}

\subsection*{You'll want to know: reservieren}
\textbf{reservieren} — \textquotedblleft{}to book, to reserve\textquotedblright{}.

\begin{quote}
Für die Nacht reserviert er ein Bett im Schlafwagen.
\end{quote}

\begin{quote}
\emph{For the night he books a bed in the sleeping car.}
\end{quote}

Past participle without \emph{ge-}: \emph{ich habe reserviert}. \emph{Die Reservierung} — the reservation.

\begin{cousinweb}
\emph{reservieren} is a learned borrowing from Latin \emph{reservāre} 'to keep back, to save up' (Lewis \& Short), probably through French \emph{réserver}.
\end{cousinweb}

\begin{grammarlens}[title={What you've built}]
One more word for telling, explaining and describing at B1.
\end{grammarlens}

\subsection*{Your turn}
\begin{itemize}
\raggedright
  \item \emph{Say it:} \emph{reservieren}
  \item \emph{Say it:} \emph{reservieren}, once more
  \item \emph{Say it:} a sentence of your own with \emph{reservieren}
  \item \emph{Recall:} say the German for the dining car, then the German for the end of the line, then say \emph{reservieren} again
\end{itemize}

\begin{tcolorbox}[breakable,colback=teal!4,colframe=teal!35!black,title={Before you move on}]
What does reservieren mean? (\textquotedblleft{}To book, to reserve\textquotedblright{}.) Say it once more.
\end{tcolorbox}
\section[weiterfahren]{weiterfahren — to travel on, to drive on}

\label{lesson:GE-C413-weiterfahren}



\begin{quote}
Before the new one: say the German for the end of the line, then the German for to book.
\end{quote}

\subsection*{You'll want to know: weiterfahren}
\textbf{weiterfahren} — \textquotedblleft{}to travel on, to drive on\textquotedblright{}.

\begin{quote}
Nach einer kurzen Pause fahren wir weiter.
\end{quote}

\begin{quote}
\emph{After a short break we drive on.}
\end{quote}

Separable; Perfekt with \emph{sein}: \emph{wir sind weitergefahren}.

\begin{cousinweb}
\emph{weiterfahren} is built from \emph{weiter} 'further, on' and the inherited verb \emph{fahren}, Proto-Germanic \emph{faraną} (EtymDB).
\end{cousinweb}

\begin{grammarlens}[title={What you've built}]
One more word for telling, explaining and describing at B1.
\end{grammarlens}

\subsection*{Your turn}
\begin{itemize}
\raggedright
  \item \emph{Say it:} \emph{weiterfahren}
  \item \emph{Say it:} \emph{weiterfahren}, once more
  \item \emph{Say it:} a sentence of your own with \emph{weiterfahren}
  \item \emph{Recall:} say the German for the end of the line, then the German for to book, then say \emph{weiterfahren} again
\end{itemize}

\begin{tcolorbox}[breakable,colback=teal!4,colframe=teal!35!black,title={Before you move on}]
What does weiterfahren mean? (\textquotedblleft{}To travel on, to drive on\textquotedblright{}.) Say it once more.
\end{tcolorbox}
\section[(dialogue)]{Words from chapters 406–413 — a review}

\label{lesson:GE-R413-review-406-413}



\begin{quote}
Say the German for to park and to pull out.
\end{quote}

\subsection*{Your turn}
Say these aloud:

\begin{itemize}
\raggedright
  \item to park, to pull out, the parking space, the multi-storey car park, the parking ticket
  \item to speed, the speed camera, to catch, the fine, the penalty ticket
  \item the right of way, to overtake, to tailgate, the oncoming traffic, to swerve
  \item the road works, the diversion, to divert, to close, the motorway service area
  \item the rear-end collision, to drive into the back, to skid, the pedestrian, the cyclist
  \item the ticket machine, the travel ticket, the fare, to stamp, to get a ticket
  \item the platform, the announcement, the compartment, the conductor, to check
  \item the sleeping car, the dining car, the end of the line, to book, to travel on
\end{itemize}

\begin{tcolorbox}[breakable,colback=teal!4,colframe=teal!35!black,title={Before you move on}]
To park? (\textbf{einparken}.) To pull out? (\textbf{ausparken}.) The parking space? (\textbf{die Parklücke}.) The multi-storey car park? (\textbf{das Parkhaus}.) The parking ticket? (\textbf{der Parkschein}.) To speed? (\textbf{rasen}.) The speed camera? (\textbf{der Blitzer}.) To catch? (\textbf{erwischen}.) The fine? (\textbf{das Bußgeld}.) The penalty ticket? (\textbf{der Strafzettel}.) The right of way? (\textbf{die Vorfahrt}.) To overtake? (\textbf{überholen}.) To tailgate? (\textbf{drängeln}.) The oncoming traffic? (\textbf{der Gegenverkehr}.) To swerve? (\textbf{ausweichen}.) The road works? (\textbf{die Baustelle}.) The diversion? (\textbf{die Umleitung}.) To divert? (\textbf{umleiten}.) To close? (\textbf{sperren}.) The motorway service area? (\textbf{die Raststätte}.) The rear-end collision? (\textbf{der Auffahrunfall}.) To drive into the back? (\textbf{auffahren}.) To skid? (\textbf{schleudern}.) The pedestrian? (\textbf{der Fußgänger}.) The cyclist? (\textbf{der Radfahrer}.) The ticket machine? (\textbf{der Fahrkartenautomat}.) The travel ticket? (\textbf{der Fahrschein}.) The fare? (\textbf{der Fahrpreis}.) To stamp? (\textbf{entwerten}.) To get a ticket? (\textbf{lösen}.) The platform? (\textbf{der Bahnsteig}.) The announcement? (\textbf{die Durchsage}.) The compartment? (\textbf{das Abteil}.) The conductor? (\textbf{der Schaffner}.) To check? (\textbf{kontrollieren}.) The sleeping car? (\textbf{der Schlafwagen}.) The dining car? (\textbf{der Speisewagen}.) The end of the line? (\textbf{die Endstation}.) To book? (\textbf{reservieren}.) To travel on? (\textbf{weiterfahren}.)
\end{tcolorbox}
